\section{Proof of the main theorem and corollaries}\label{sec:proofs}

Throughout this section $j=10(B+\ell+100)$ is the integer of \cref{thm:halting}, with $\ell=\binom{276}{2}+1=37951$, and
\[
k=\max(2,2j)=2j,\qquad q=k+1,\qquad a=2q^2+1,\qquad D_0=50q,\qquad \eps_0=\eps_k=(4D_0^2)^{-a}.
\]
The dimension is $d=2+k$. The family $M\mapsto R_M$ is the one of \cref{thm:width}. We write $\delta_O(t)$ as in \cref{set:def-forget}, adding a superscript $M$ when the Wang set is $R_M$.

\begin{lemma}[One threshold]\label{pf:threshold}
Let $0<\eps\le\eps_0$. Then $\beta=2D_0^2\eps^{1/a}\le\frac12$, the hypothesis of \cref{thm:memory} holds, and $\eps<10^{-9}$. Both conditions on the rate have the form $\eps\le\eps_0$, so they hold at every rate in $(0,\eps_0]$.
\end{lemma}

\begin{proof}
The map $\eps\mapsto\beta$ is increasing, and at $\eps_0$ it equals $2D_0^2(4D_0^2)^{-1}=\frac12$. \Cref{thm:memory} asks for $0<\eps\le(4D_0^2)^{-a}=\eps_0$. Since $q\ge3$, $D_0\ge150$ and $a\ge19$, we get $\eps_0\le(90000)^{-19}<10^{-9}$.
\end{proof}

\begin{theorem}[Theorem A]\label{thm:main}
For every $\eps\in(0,\eps_0]$ and every Turing machine $M$, the noisy quorum rule $F^{[k]}_{R_M}$ on $\Z^{2+k}$ with uniform replacement at rate $\eps$ is a positive-rate probabilistic cellular automaton, and it is ergodic if and only if $M$ halts. The map $M\mapsto F^{[k]}_{R_M}$ is total computable.
\end{theorem}

\begin{proof}[Proof of \cref{thm:main}]
Fix $\eps\in(0,\eps_0]$ and use uniform replacement noise at rate $\eps$. The map $M\mapsto R_M$ is total computable by \cref{thm:width}, and $T_R$, $\A_R$ and the local rule of $F^{[k]}_R$ on $\mathcal N=\Lambda\times\{c_0,\dots,c_k\}$ are computed from $R$ by finite enumeration. So $M\mapsto F^{[k]}_{R_M}$ is total computable.

Suppose $M$ halts. By \cref{thm:width}, $R_M$ is $j$-token refutable; let $\tau=\tau(R_M)$ be its refutation time. The hypotheses of \cref{thm:forgetting} hold: $j\ge1$, $k\ge\max(2,2j)$, $0<\eps<1$, and $\beta\le\frac12<1$ by \cref{pf:threshold}. So for every window $O$ and every $t\ge0$,
\begin{equation}\label{pf:eq-forget}
\delta^M_O(t)\le2|O|\,\eps^{-\tau/a}\,\frac{\beta^t}{1-\beta}\le4|O|\,\eps^{-\tau/a}\,2^{-t},
\end{equation}
which tends to $0$. The process forgets, and it is ergodic by \cref{set:forget-ergodic}.

Suppose $M$ does not halt. By \cref{thm:width}, $R_M$ tiles; let $\xi$ be the cloned tiling configuration of a tiling. By \cref{thm:memory} and \cref{pf:threshold}, for every site $v$ and every $t\ge0$,
\begin{equation}\label{pf:eq-memory}
\Prob_\xi\bigl(X_t(v)\ne\xi(v)\bigr)<1.01\,\eps,\qquad \Prob_\bot\bigl(X_t(v)\text{ is live}\bigr)<1.01\,\eps .
\end{equation}
Take the window $\{v\}$ and the event that $v$ carries the live symbol $\xi(v)$. Then $\Prob_\xi(X_t(v)=\xi(v))>1-1.01\eps$ and $\Prob_\bot(X_t(v)=\xi(v))<1.01\eps$, so the process has memory with $\alpha=1-2.02\eps>0$. By \cref{set:memory-nonergodic} it has two different invariant measures, and it is not ergodic.

Both cases hold for every $\eps\in(0,\eps_0]$, and $j$, $k$ and $\eps_0$ do not depend on $M$.
\end{proof}

\begin{remark}[Scope of the constants]\label{pf:scope}
Since $B\ge1$, we have $j\ge10(1+37951+100)=380520$, so the dimension $d=2+2j$ given by \cref{thm:halting} is a fixed integer, at least $761042$. Its exact value depends on the constant $B$ of \cref{thm:parser}, a fixed integer, shown to exist, that we do not evaluate. With the three-token refutation of \cref{sec:three}, the same family gives $j=3$ and $d=8$ (\cref{cor:eight}). The rate $\eps_0$ depends only on $k$, and $\beta$ only on $k$ and $\eps$; the Wang set enters the forgetting bound only through $\tau(R_M)$. The alphabet $\A_{R_M}$ varies with $M$, so the least transition probability $\eps/|\A_{R_M}|$ is not bounded below over the family. A fixed noise rate, as in the title, means a fixed replacement rate $\eps$ with uniform replacement on the alphabet of each member.
\end{remark}

\subsection{Corollaries}\label{pf:sec-cor}

\begin{corollary}[Every small rate]\label{cor:rates}
For every machine $M$ the following are equivalent:
\begin{itemize}
\item[(i)] $M$ halts;
\item[(ii)] $F^{[k]}_{R_M}$ with uniform replacement noise is ergodic at some rate $\eps\in(0,\eps_0]$;
\item[(iii)] it is ergodic at every rate $\eps\in(0,\eps_0]$.
\end{itemize}
Consequently, for each fixed $\eps\in(0,\eps_0]$, no algorithm decides, given a finite Wang set $R$, whether $F^{[k]}_R$ with uniform replacement noise at rate $\eps$ is ergodic.
\end{corollary}

\begin{proof}
By \cref{pf:threshold} the hypotheses of \cref{thm:forgetting,thm:memory} used in the proof of \cref{thm:main} hold at every rate in $(0,\eps_0]$. So (i) implies (iii), (iii) implies (ii), and if $M$ does not halt the process is not ergodic at any rate in $(0,\eps_0]$, which is (ii) implies (i). For the last claim, a decision procedure at rate $\eps$, composed with the computable map $M\mapsto R_M$, would decide halting.
\end{proof}

\begin{corollary}[Every Wang set at once]\label{cor:tiling}
There is a total computable map $R\mapsto G_R$ from finite Wang sets to positive-rate probabilistic cellular automata on $\Z^{2+k}$, each with uniform replacement noise, such that for every $\eps\in(0,\eps_0]$ and every finite Wang set $R$, the automaton $G_R$ at rate $\eps$ is ergodic if and only if $R$ does not tile the plane.
\end{corollary}

\begin{proof}
Let $M_R$ be the machine that, for $n=1,2,\dots$ in turn, checks whether $R$ has a legal pattern on the $n\times n$ square, and halts at the first $n$ for which it has none. The map $R\mapsto M_R$ is computable, and each check is finite. If $R$ tiles, restrictions of a tiling give legal squares of every size, so $M_R$ does not halt. If $M_R$ does not halt, $R$ has legal squares of every size; by compactness of $R^{\Z^2}$ (König's lemma on the finitely branching tree of legal centred squares) it tiles. So $M_R$ halts if and only if $R$ does not tile. Put $G_R=F^{[k]}_{R_{M_R}}$ with uniform replacement noise, and apply \cref{thm:main} to $M_R$.
\end{proof}

The automaton $G_R$ is built from the recoded set $R_{M_R}$, not from $R$ itself. So \cref{cor:tiling} gives one fixed dimension for every Wang set, whatever its token width, through the family of \cref{thm:width}; it says nothing new about the quorum rule of $R$ itself.

The next corollary concerns other noise. For a probability law $\nu$ on $\A_R$, \emph{$\nu$-replacement noise} at rate $\eps$ is defined as in \cref{set:def-noise}, with the values $\zeta_t(v)$ independent with law $\nu$ in place of the uniform law. Uniform replacement is the case where $\nu$ is uniform. \Cref{set:feller,set:forget-ergodic,set:memory-nonergodic} and their proofs hold verbatim for $\nu$-replacement noise: they use only the mixture formula and the fact that $X_1|_O$ depends on finitely many coordinates of $X_0$ and on independent noise.

\begin{corollary}[Other replacement laws]\label{cor:laws}
For each machine $M$ let $\nu_M$ be a probability law on $\A_{R_M}$, and let $\eps\in(0,\eps_0]$.
\begin{itemize}
\item[(i)] $F^{[k]}_{R_M}$ with $\nu_M$-replacement noise at rate $\eps$ is ergodic if and only if $M$ halts. Every transition probability of this process is at least $\eps\min_s\nu_M(s)$, which is positive when $\nu_M$ has full support.
\item[(ii)] Suppose $M$ does not halt. Let the replacement events be as in \cref{set:def-noise}, and let the values written at events be arbitrary random variables, which may depend on the event indicators, on the past of the run and on further randomness. Then the bounds \eqref{pf:eq-memory} hold for the run from every cloned tiling configuration $\xi$ and for the run from $\bot$, at every site and time.
\end{itemize}
\end{corollary}

\begin{proof}
(ii) The proof of \cref{thm:memory} has a deterministic part and a counting part. The deterministic part (\cref{mem:two-parents}) holds for every realization of the events and of the values: a bad point that is not an event has two bad points one step earlier, in two different provider planes, and no point of time $0$ is bad. The counting part uses the noise only through the bound $\Prob(\text{every point of }S\text{ is an event})\le\eps^{|S|}$ for finite sets $S$ of points of positive time (\cref{mem:laws}). Both parts hold under the stated hypotheses, so the bounds hold.

(i) Suppose $M$ does not halt. Values drawn independently from $\nu_M$ are a special case of (ii), so \eqref{pf:eq-memory} holds, the process has memory as in the proof of \cref{thm:main}, and \cref{set:memory-nonergodic} shows that it is not ergodic.

Suppose $M$ halts. In the proof of \cref{thm:forgetting}, the runs from $x$ and from $\bot$ use the same indicators $\eta$ and the same values $\zeta$. The covering field $E$ is a function of $\eta$ alone (\cref{def:domination}). The transport and cover lemmas and \cref{fgt:live} hold for every realization of $\eta$ and $\zeta$. The discrepancy path of \cref{fgt:path} uses only that both runs take the common value at an event. The final count uses only the independence of the indicators and $\Prob(\eta_t(v)=1)=\eps$. None of these steps uses the law of $\zeta$. So \eqref{fgt:eq-main} and \eqref{fgt:eq-tv} hold for $\nu_M$-replacement noise with the same constants, the process forgets, and \cref{set:forget-ergodic} shows that it is ergodic.

Finally, given $X_t$, the value $X_{t+1}(v)$ equals $s$ with probability $(1-\eps)\mathbf 1[F^{[k]}_{R_M}(X_t)(v)=s]+\eps\,\nu_M(s)\ge\eps\,\nu_M(s)$.
\end{proof}

The forgetting half needs the values to be drawn without looking at the configuration, because its coupling gives both runs the same value at each event. The memory half needs nothing about the values.

\begin{corollary}[Mixing times are not computably bounded]\label{cor:mixing}
Let $\eps\in(0,\eps_0]$ be rational, and use uniform replacement noise at rate $\eps$.
\begin{itemize}
\item[(i)] For every machine $M$, window $O$ and time $t$, the supremum in $\delta^M_O(t)$ is a maximum over finitely many restrictions of the start, and $\delta^M_O(t)$ is a rational number computable from $M$, $O$ and $t$.
\item[(ii)] Let $t_{\mathrm{mix}}(M)$ be the least $t_0\ge0$ such that $\delta^M_{\{0\}}(t)<\frac12$ for every $t\ge t_0$, and $t_{\mathrm{mix}}(M)=\infty$ if there is none. Then $t_{\mathrm{mix}}(M)<\infty$ if and only if $M$ halts.
\item[(iii)] No computable function $b$ satisfies $t_{\mathrm{mix}}(M)\le b(M)$ for every halting $M$. No computable function $b$ satisfies $\tau(R_M)\le b(M)$ for every halting $M$.
\end{itemize}
\end{corollary}

\begin{proof}
(i) Let $\mathcal N=\Lambda\times\{c_0,\dots,c_k\}$ and $\mathcal N_t$ be as in \cref{sec:setting}, and write $\A=\A_{R_M}$ and $F=F^{[k]}_{R_M}$. Given $X_s$, the values $X_{s+1}(v)$ are independent, and $X_{s+1}(v)=s'$ with probability $(1-\eps)\mathbf 1[F(X_s)(v)=s']+\eps/|\A|$, where $F(X_s)(v)$ depends on $X_s$ only on $v+\mathcal N$. Hence, for a window $S$ and $w\in\A^S$,
\[
\Prob_x\bigl(X_{s+1}|_S=w\bigr)=\sum_{y\in\A^{S+\mathcal N}}\Prob_x\bigl(X_s|_{S+\mathcal N}=y\bigr)\prod_{v\in S}\Bigl((1-\eps)\mathbf 1\bigl[F(y)(v)=w(v)\bigr]+\frac{\eps}{|\A|}\Bigr).
\]
where $F(y)(v)$ is the local rule applied to $y$ on $v+\mathcal N$. Apply this $t$ times, with $S=O+\mathcal N_{t-s-1}$ at step $s$, from $X_0|_{O+\mathcal N_t}=x|_{O+\mathcal N_t}$ to $X_t|_O$. Each step is a finite sum of products of rationals computable from $M$. So $\mathrm{Law}_x(X_t|_O)$ is a rational vector, computable from $M$, $O$, $t$ and $x|_{O+\mathcal N_t}$, and it depends on $x$ only through this restriction. The supremum over $x$ is a maximum over the finite set $\A^{O+\mathcal N_t}$, and the total variation distance of two rational vectors is rational.

(ii) If $M$ halts, \eqref{pf:eq-forget} with $O=\{0\}$ gives $\delta^M_{\{0\}}(t)\to0$, so $t_{\mathrm{mix}}(M)<\infty$. If $M$ does not halt, let $\xi$ be a cloned tiling configuration. By \eqref{pf:eq-memory}, for every $t$,
\[
\delta^M_{\{0\}}(t)\ge\Prob_\xi\bigl(X_t(0)=\xi(0)\bigr)-\Prob_\bot\bigl(X_t(0)=\xi(0)\bigr)>1-2.02\,\eps>\tfrac12,
\]
so $t_{\mathrm{mix}}(M)=\infty$.

(iii) Suppose $b$ is computable and $t_{\mathrm{mix}}(M)\le b(M)$ for every halting $M$. Given $M$, compute $\delta^M_{\{0\}}(b(M))$ by (i) and answer ``halts'' exactly when it is below $\frac12$. If $M$ halts, then $b(M)\ge t_{\mathrm{mix}}(M)$, so the value is below $\frac12$. If $M$ does not halt, the value exceeds $\frac12$ by the proof of (ii). This decides the halting problem, which is impossible.

For the refutation time, let $M$ halt and $\tau=\tau(R_M)$. By \eqref{pf:eq-forget} with $|O|=1$, $\delta^M_{\{0\}}(t)<\frac12$ whenever $2^t>8\eps^{-\tau/a}$, that is, whenever $t>3+(\tau/a)\log_2(1/\eps)$. So $t_{\mathrm{mix}}(M)\le4+\tau L$ with $L=\lceil\log_2(1/\eps)\rceil$, since $a\ge1$. A computable $b$ with $\tau(R_M)\le b(M)$ for every halting $M$ would give the computable bound $4+b(M)L$ on $t_{\mathrm{mix}}(M)$, which we have excluded.
\end{proof}

The process and the all-bottom start are invariant under translations of $\Z^{2+k}$, so $\delta^M_{\{v\}}=\delta^M_{\{0\}}$ for every site $v$. Restricting to a smaller window does not increase total variation, so $\delta^M_O\ge\delta^M_{\{0\}}$ for every nonempty window $O$. The conclusions of \cref{cor:mixing} therefore hold for the mixing time of every nonempty window. Forgetting happens for every halting member, but its onset is not computably bounded. The bound \eqref{pf:eq-forget} reads $\delta^M_O(t)\le2^{-(t-t_1)}$ with $t_1=(\tau(R_M)/a)\log_2(1/\eps)+\log_2(4|O|)$: a delay set by the refutation time, followed by decay at a rate that does not depend on $M$.

Ergodicity also has an equilibrium reading. A PCA on $\Z^d$ with a strictly positive transition table defines a Gibbs specification on $\Z^{d+1}$, its space-time interaction, in which time is one more lattice direction \cite{DM26}. For $F^{[k]}_{R_M}$ with uniform replacement noise at rate $\eps$, the table is
\[
p_M(s\mid y)=(1-\eps)\,\mathbf 1\bigl[s=f_M(y)\bigr]+\frac{\eps}{|\A_{R_M}|},\qquad s\in\A_{R_M},\ y\in\A_{R_M}^{\mathcal N},
\]
where $f_M$ is the local rule of $F^{[k]}_{R_M}$ on $\mathcal N$. Each entry is at least $\eps/|\A_{R_M}|>0$. Write a space-time point as $(v,t)\in\Z^{2+k}\times\Z$; its parents are the points $(v+w,t-1)$, $w\in\mathcal N$. The \emph{space-time interaction} $\Psi_M$ assigns to each point $(v,t)$ the energy $-\log p_M\bigl(\omega(v,t)\mid\omega\text{ at the parents of }(v,t)\bigr)$ of a configuration $\omega\in\A_{R_M}^{\Z^{3+k}}$. It has finite range and is invariant under all translations of $\Z^{3+k}$. A \emph{Gibbs measure} for $\Psi_M$ is a DLR state of its specification, as defined in \cite{DM26}.

\begin{corollary}[Space-time Gibbs measures]\label{cor:gibbs}
Let $\eps\in(0,\eps_0]$. If $M$ halts, $\Psi_M$ has exactly one translation-invariant Gibbs measure. If $M$ does not halt, it has at least two. For rational $\eps$, the table $p_M$ is computable from $M$; so it is undecidable whether a member of this computable family of finite-range, translation-invariant interactions on $\Z^{3+k}$ has a unique translation-invariant Gibbs measure.
\end{corollary}

\begin{proof}
We use three results of \cite{DM26}, stated there for every spatial dimension, every finite alphabet with at least two symbols, every finite neighbourhood containing the origin, and every strictly positive table. Our process fits: the alphabet is finite, $\mathcal N$ is finite and contains $0=(0,c_0)$, and every entry of $p_M$ is positive. The three results are: the lemma that PCA histories satisfy the DLR equations, so that every two-sided stationary trajectory law is a Gibbs measure; the corollary that periodic states are stationary trajectories, which with that lemma identifies the translation-invariant Gibbs measures with the trajectory laws of translation-invariant stationary laws, the correspondence of Goldstein, Kuik, Lebowitz and Maes \cite{GKLM89}; and item~(i) of the theorem on the Gibbs simplex of an ergodic PCA: an ergodic strictly positive PCA has exactly one periodic, in particular one translation-invariant, Gibbs measure. The notion of ergodicity in \cite{DM26}, a stationary law that attracts every initial law, is the one of \cref{set:def-ergodic}.

If $\A_{R_M}=\{\bot\}$, then $R_M$ has no legal $3\times3$ block, so it does not tile and $M$ halts by \cref{thm:width}. The configuration space is then one point, which carries exactly one measure, and the claim holds. Otherwise $\A_{R_M}$ has at least three symbols.

Suppose $M$ halts. By \cref{thm:main} the process is ergodic, with invariant measure $\pi$. The process commutes with the translations of $\Z^{2+k}$, so every translate of $\pi$ is invariant, and $\pi$ is translation invariant by uniqueness. Its two-sided stationary trajectory law is invariant under all translations of $\Z^{3+k}$ and is a Gibbs measure. By the theorem on the Gibbs simplex, item~(i), it is the only translation-invariant one.

Suppose $M$ does not halt. By \cref{thm:width}, $R_M$ tiles, so the set $\mathcal T_M\subseteq R_M^{\Z^2}$ of tilings is nonempty. It is closed in the compact product space and invariant under the shifts of $\Z^2$. Fix $z_0\in\mathcal T_M$ and let $m_n$ be the average of the point masses at the translates $z_0(\cdot+w)$, $w\in[-n,n]^2$. A weak limit point $m$ of $(m_n)$ is a probability measure on $\mathcal T_M$, and it is shift invariant: for a unit vector $e$ of $\Z^2$, the measures $m_n$ and their translates by $e$ differ by at most $2(2n+1)/(2n+1)^2$ in total variation. For a tiling $z$ let $\xi_z$ be its cloned tiling configuration. The map $z\mapsto\xi_z$ is continuous, since $\xi_z(u,P)$ depends on $z$ only on $u+\Lambda$. It satisfies $\xi_{z(\cdot+w)}=\xi_z(\cdot+(w,0))$, and $\xi_z$ is invariant under translations of the transverse coordinate. So the image $\kappa$ of $m$ is a probability measure on $\Omega$ that is invariant under all translations of $\Z^{2+k}$.

Fix a site $v$. The function $x\mapsto\Prob_x(X_t(v)\text{ is live})$ depends on finitely many coordinates of $x$, and by \eqref{pf:eq-memory} it exceeds $1-1.01\eps$ at every $\xi_z$. By the mixture formula of \cref{sec:setting}, $\Prob_\kappa(X_t(v)\text{ is live})\ge1-1.01\eps$ for every $t$. The laws $\kappa\Phi^t$ are translation invariant, because $\kappa$ is and the process commutes with translations. By \cref{set:feller}, the averages $\frac1n\sum_{t<n}\kappa\Phi^t$ have a weak limit point $\pi^+$, which is invariant. It is translation invariant, as a weak limit of translation-invariant measures, and $\pi^+(\omega(v)\text{ is live})\ge1-1.01\eps$, because the cylinder is open and closed. In the same way the averages from the all-bottom configuration, which is translation invariant, have an invariant, translation-invariant limit point $\pi^-$, with $\pi^-(\omega(v)\text{ is live})\le1.01\eps$ by \eqref{pf:eq-memory}. As $1.01\eps<1-1.01\eps$, the measures $\pi^+$ and $\pi^-$ differ.

For an invariant $\pi$, its two-sided stationary trajectory law is the measure on $\A_{R_M}^{\Z^{3+k}}$ whose restriction to the times $t_0,\dots,t_0+n$ is the law of $(X_0,\dots,X_n)$ under $\mathrm{Law}_\pi$; these restrictions are consistent because $\pi$ is invariant, and Kolmogorov's theorem gives the measure. For $\pi=\pi^\pm$ it is invariant under time translations, and under space translations because $\pi^\pm$ is and the process commutes with them. So both trajectory laws are translation-invariant Gibbs measures. Their marginals at time $0$ are $\pi^+\ne\pi^-$, so they differ.

The last claim follows: an algorithm deciding uniqueness of the translation-invariant Gibbs measure for the tables $p_M$ would decide halting.
\end{proof}

\begin{remark}[Arithmetical placement]\label{pf:arith}
For a PCA with a strictly positive rational table, ergodicity is a $\Pi^0_2$ property of the table \cite{DM26}. The reason is short. By Halawi's theorem \cite{Halawi22}, an ergodic positive PCA is uniformly ergodic. So ergodicity says: for every window $O$ and every $m\ge1$ there is $t$ with $\Delta_O(t)<1/m$, where $\Delta_O(t)=\sup_{x,y}\dtv\bigl(\mathrm{Law}_x(X_t|_O),\mathrm{Law}_y(X_t|_O)\bigr)$. The argument of \cref{cor:mixing}(i) computes $\Delta_O(t)$ exactly, and $\Delta_O(t)$ does not increase with $t$ \cite{DM26}. \Cref{thm:main} gives a lower bound at the first level. For each rational $\eps\in(0,\eps_0]$, the computable map $M\mapsto F^{[k]}_{R_M}$ reduces the halting problem to ergodicity at rate $\eps$. So ergodicity of positive-rate PCA on $\Z^{2+k}$ at the fixed rate $\eps$, with uniform replacement, is $\Sigma^0_1$-hard, and non-ergodicity is $\Pi^0_1$-hard; in particular the non-ergodic members do not form a computably enumerable set. Whether ergodicity at a fixed rate is $\Pi^0_2$-complete is not settled by \cref{thm:main}.
\end{remark}

\subsection{Perturbed kernels}\label{pf:sec-robust}

Theorem~A fixes one noise, uniform replacement at rate $\eps$. The dichotomy does
not depend on that choice. It holds for every probabilistic cellular automaton whose
table lies within a fixed total-variation distance of the noiseless rule, with a
radius that depends on $k$ only. The proof reuses the proofs of
\cref{thm:memory,thm:forgetting}. The memory half needs nothing beyond
\cref{mem:laws}. The forgetting half needs two changes: the discrepancy path of
\cref{fgt:path} may pass through faults, and a fault on that path receives a contour
with no edges.

Throughout this subsection $R$ is a Wang set, $k\ge2$, $q=k+1$, $a=2q^2+1$,
$D_0=50q$, $\A=\A_R$, $F=F^{[k]}_R$ and $\mathcal N=\Lambda\times\{c_0,\dots,c_k\}$.
For $x\in\Omega$ and a site $v$ let $x_v\in\A^{\mathcal N}$ be given by
$x_v(w)=x(v+w)$, and let $f\colon\A^{\mathcal N}\to\A$ be the local rule of $F$, so
that $F(x)(v)=f(x_v)$.

\begin{definition}[Local map fields]\label{pf:def-field}
A \emph{local map field} consists of random variables $\chi_t(v)\in\{0,1\}$ and
random maps $f_{t,v}\colon\A^{\mathcal N}\to\A$, for $t\ge1$ and
$v\in\Z^{2+k}$, all defined on one probability space, such that $f_{t,v}=f$
whenever $\chi_t(v)=0$. A point $(v,t)$ with $\chi_t(v)=1$ is a \emph{fault}. The
\emph{run} from $x\in\Omega$ is $X_0=x$ and $X_t(v)=f_{t,v}\bigl((X_{t-1})_v\bigr)$
for $t\ge1$. Runs from several starts driven by one field form the
\emph{common-map coupling}. For $0<\gamma<1$ the field is \emph{$\gamma$-bounded} if
\[
\Prob\bigl(\chi_t(v)=1\text{ for every }(v,t)\in S\bigr)\le\gamma^{|S|}
\]
for every finite set $S$ of points of positive time.
\end{definition}

Nothing else is assumed. Faults may be correlated in space and in time, and the map
at a fault may be any function of its inputs, chosen at random jointly with all
other maps. Bounds of this form are the noise condition of Toom's stability theorem \cite{Toom80}; see also \cite{Gacs21}. Replacement noise at rate $\eps$ (\cref{set:def-noise}) is the
$\eps$-bounded field with $\chi_t(v)=\eta_t(v)$, in which $f_{t,v}$ is the constant
map with value $\zeta_t(v)$ when $\eta_t(v)=1$.

\begin{lemma}[Common local maps]\label{pf:field}
Let $0<\gamma<1$, let $(\chi,f_{t,v})$ be a $\gamma$-bounded local map field, and
let $E$ be the domination field of \cref{def:domination} computed from $\chi$ in
place of $\eta$.
\begin{enumerate}[label=\textup{(\roman*)}]
\item Let $R$ tile, let $\xi$ be a cloned tiling configuration, and suppose
$\lambda=2D_0^2\gamma^{1/q^2}<1$. Let $X$ be the run from $\xi$ and $Y$ the run from
$\bot$. For every site $v$ and every $t\ge0$,
\[
\Prob\bigl(X_t(v)\ne\xi(v)\bigr)\le\gamma\Bigl(1+\frac{2\lambda}{1-\lambda}\Bigr),
\qquad
\Prob\bigl(Y_t(v)\ne\bot\bigr)\le\gamma\Bigl(1+\frac{2\lambda}{1-\lambda}\Bigr).
\]
\item Let $j\ge1$, $k\ge\max(2,2j)$, let $R$ be $j$-token refutable with time
$\tau$, and suppose $\beta=2D_0^2\gamma^{1/a}<1$. Let $X$ and $Y$ be the runs from
configurations $x$ and $y$ under the common-map coupling. For every site $v$ and
every $t\ge0$,
\[
\Prob\bigl(X_t(v)\ne Y_t(v)\bigr)\le2\gamma^{-\tau/a}\,\frac{\beta^t}{1-\beta}.
\]
\end{enumerate}
\end{lemma}

\begin{proof}
In both parts, read ``event'' as ``fault'' in \cref{sec:memory,sec:forgetting}.
At a point that is not a fault the run applies $F$:
$X_t(v)=f\bigl((X_{t-1})_v\bigr)=F(X_{t-1})(v)$.

(i) The proof of \cref{mem:two-parents} uses a point that is not an event only
through $X_t(u,P)=F(X_{t-1})(u,P)$, and the start only at time $0$. It never uses
the value written at an event. So \cref{mem:two-parents} and \cref{mem:dominate}
hold for every realization of the field. The proof of \cref{mem:field-bound} uses
the indicators only through the recursive definition of $E$ and through the bound
that every finite set $S$ of points of positive time consists of events with
probability at most $\eps^{|S|}$ (\cref{mem:laws}). Here that bound holds with $\gamma$ in
place of $\eps$. With $|S|=1$ it also gives $\Prob(\chi_t(v)=1)\le\gamma$. So the
proof of \cref{mem:field-bound} gives
$\Prob(E_t(v)=1)\le\gamma(1+2\lambda/(1-\lambda))$, and \cref{mem:dominate} gives (i).

(ii) If $L_R=\emptyset$, then $\A=\{\bot\}$ and there is nothing to prove; otherwise
$\tau\ge1$, as in the proof of \cref{thm:forgetting}. Fix $r$ with
$S_\tau(0,\dots,0)=\emptyset$ for the box $V=\{-r,\dots,r\}^2$.

\emph{Step 1: the cover.} \Cref{fgt:transport}, \cref{fgt:cover} and \cref{fgt:live}
are statements about one run and one realization. Their proofs use a point that is
not an event only through $X_t(v)=F(X_{t-1})(v)$ (Step~2 of the proof of
\cref{fgt:transport}), and an event only through $E=1$ there. They never use the
value written at an event. So they hold for every realization of the field and for
the run from every start, with the one field $E$, which depends on $\chi$ alone.
Hence, for $t\ge\tau$, every live value $X_t(v)$ or $Y_t(v)$ has $E_t(v)=1$.

\emph{Step 2: the discrepancy path.} Let $t\ge1$ and $X_t(v)\ne Y_t(v)$. Since
$X_t(v)=f_{t,v}\bigl((X_{t-1})_v\bigr)$ and $Y_t(v)=f_{t,v}\bigl((Y_{t-1})_v\bigr)$
with one and the same map $f_{t,v}$, the inputs differ:
$X_{t-1}(v')\ne Y_{t-1}(v')$ for some $v'\in v+\mathcal N=\mathcal N(v)$. This holds
whether or not $(v,t)$ is a fault. Repeating down to time $0$ gives sites
$v_t=v,v_{t-1},\dots,v_0$ with $v_{s-1}\in\mathcal N(v_s)$ and
$X_s(v_s)\ne Y_s(v_s)$ for $0\le s\le t$. So $(v_s,s)$ and $(v_{s-1},s-1)$ are
joined in $\mathcal G$. This replaces \cref{fgt:path}. The points $(v_s,s)$ may now
be faults.

\emph{Step 3: marks on the path.} Let $\tau\le s\le t$. If $(v_s,s)$ is a fault,
then $E_s(v_s)=1$ by \cref{def:domination}. Otherwise, $X_s(v_s)$ and $Y_s(v_s)$ are
different symbols of $L_R\cup\{\bot\}$, so one of them is live, and Step~1 gives
$E_s(v_s)=1$.

\emph{Step 4: packing.} Run the proof of \cref{fgt:packing} with one change. Put
$t_1=t$. Given $t_l\ge\tau$ and $\rho_l=(v_{t_l},t_l)$:
\begin{itemize}
\item If $\rho_l$ is not a fault, apply \cref{lem:contour}, with $\mathcal U$ the
marked points and $\mathcal U_{\mathrm f}$ the faults, at $\rho_l$. This gives
$m_l$, $\Sigma_l$, $K_l$ and $s_l$ as there.
\item If $\rho_l$ is a fault, put $K_l=\emptyset$, $\Sigma_l=\{\rho_l\}$, $m_l=1$
and $s_l=t_l$.
\end{itemize}
In both cases:
\begin{itemize}
\item $\Sigma_l$ is a set of $m_l$ faults with times in $[s_l,t_l]$, and every point
of $\Sigma_l$ is $\rho_l$ or a vertex of $K_l$;
\item every vertex of $K_l$ has time in $[s_l,t_l]$;
\item if $K_l\ne\emptyset$, its edges form a connected graph with $\rho_l$ among its
vertices (\cref{lem:contour}(i));
\item and
\[
|K_l|\le q^2(m_l-1),\qquad t_l-s_l\le q^2(m_l-1).
\]
\end{itemize}
For the second case both sides of each inequality are $0$. These are the only
properties of the contours that the proof of \cref{fgt:packing} uses. Continue with
$t_{l+1}=s_l-1$ while $s_l-1\ge\tau$. The slabs $[s_l,t_l]$ are pairwise disjoint,
so the sets $\Sigma_l$ are pairwise disjoint. Moreover
$t-\tau+1\le\sum_l(t_l-s_l+1)\le q^2m$ with $m=\sum_l m_l$. Let $K$ consist of the
$t$ path edges together with all $K_l$. The path edges form a connected graph whose
vertices include every $\rho_l$, since $t\ge\tau\ge1$. Each nonempty $K_l$ is
connected and contains $\rho_l$. So $K$ is connected, its vertices include $(v,t)$
and $\Sigma=\bigcup_l\Sigma_l$, and $e=|K|$ satisfies $t\le e\le\tau-1+2q^2m$. Hence
$|\Sigma|=m\ge(e-\tau)/a$. The proof of \cref{fgt:packing} used the start $\bot$
only to know that one of two different values is live, and Step~3 supplies this for
any two runs.

\emph{Step 5: counting.} For $t<\tau$ the right side is at least $2$, as in the proof
of \cref{thm:forgetting}. For $t\ge\tau$, by Step~4 the event $X_t(v)\ne Y_t(v)$ is
contained in the union, over $e\ge t$, over connected sets $K$ of $e$ edges of
$\mathcal G$ with $(v,t)$ among their vertices, and over sets $\Sigma$ of at least
$(e-\tau)/a$ of their vertices, of the events $\{\chi=1\text{ on }\Sigma\}$. These
$(K,\Sigma)$ form a countable family of fixed sets. Each event has probability at
most $\gamma^{|\Sigma|}\le\gamma^{(e-\tau)/a}$, by $\gamma$-boundedness; sets
$\Sigma$ that contain a point of time $0$ have probability $0$. By \cref{fgt:count}
there are at most $D_0^{2e}2^{e+1}$ pairs $(K,\Sigma)$ for each $e$. Summing as in the
proof of \cref{thm:forgetting} gives the bound. No independence is used, only the
union bound.
\end{proof}

A \emph{local kernel} on $\A$ is a family $Q=(Q(\cdot\mid y))_{y\in\A^{\mathcal N}}$
of probability laws on $\A$. The \emph{probabilistic cellular automaton with kernel
$Q$} is the process with $X_0=x$ in which, given $X_0,\dots,X_t$, the values
$X_{t+1}(v)$ are independent and $X_{t+1}(v)=s$ with probability
$Q\bigl(s\mid(X_t)_v\bigr)$. We write $\mathrm{Law}^Q_x$, $\Prob^Q_x$ and
$\delta^Q_O(t)$ as in \cref{sec:setting}. \Cref{set:feller,set:forget-ergodic,set:memory-nonergodic}
and their proofs hold verbatim for it: they use only the mixture formula and the fact
that $X_1|_O$ depends on $X_0$ only on $O+\mathcal N$ and on independent randomness.
Uniform replacement at rate $\eps$ is the kernel
$p_\eps(s\mid y)=(1-\eps)\mathbf 1[s=f(y)]+\eps/|\A|$. On local kernels over $\A$ put
\[
d(Q,Q')=\max_{y\in\A^{\mathcal N}}\dtv\bigl(Q(\cdot\mid y),Q'(\cdot\mid y)\bigr).
\]
Then $d(Q,\delta_f)=\max_y\bigl(1-Q(f(y)\mid y)\bigr)$, where $\delta_f$ is the
noiseless rule $\delta_f(\cdot\mid y)=\delta_{f(y)}$.

\begin{corollary}[Perturbed kernels]\label{cor:robust}
Let $j$, $k$, $q$, $a$, $D_0$ and $\eps_0$ be as in this section, or as in
\cref{cor:eight}. For a machine $M$ write $\A_M=\A_{R_M}$, let $f_M$ be the local
rule of $F^{[k]}_{R_M}$, and let $p_{M,\eps}$ be uniform replacement at rate $\eps$
on $\A_M$.
\begin{enumerate}[label=\textup{(\roman*)}]
\item Let $0<\gamma\le2\eps_0$ and let $Q$ be a local kernel on $\A_M$ with
\[
Q\bigl(f_M(y)\mid y\bigr)\ \ge\ 1-\gamma\qquad\text{for every }y\in\A_M^{\mathcal N}.
\]
The probabilistic cellular automaton with kernel $Q$ is ergodic if and only if $M$
halts. More precisely, put $\beta=2D_0^2\gamma^{1/a}$; then $\beta<0.52$. If $M$
halts and $\tau=\tau(R_M)$, then for every window $O$ and every $t\ge0$
\[
\delta^Q_O(t)\le2|O|\,\gamma^{-\tau/a}\,\frac{\beta^t}{1-\beta}.
\]
If $M$ does not halt and $\xi$ is a cloned tiling configuration, then for every site
$v$ and every $t\ge0$
\[
\Prob^Q_\xi\bigl(X_t(v)\ne\xi(v)\bigr)<1.01\,\gamma,\qquad
\Prob^Q_\bot\bigl(X_t(v)\text{ is live}\bigr)<1.01\,\gamma ,
\]
and the process has memory with $\alpha=1-2.02\gamma$.
\item Let $\eps\in(0,\eps_0]$. Every local kernel $Q$ on $\A_M$ with
$d(Q,p_{M,\eps})\le\eps_0$ satisfies the hypothesis of (i) with $\gamma=2\eps_0$. So
the probabilistic cellular automaton with kernel $Q$ is ergodic if and only if $M$
halts. The radius $\eps_0$ depends on $k$ only. It is the same for every machine $M$
and every rate $\eps\in(0,\eps_0]$.
\item Let $M\mapsto Q_M$ be a total computable map to local kernels on $\A_M$ with
rational entries, each satisfying the hypothesis of (i) with one fixed
$\gamma\in(0,2\eps_0]$. Then $Q_M$ is ergodic if and only if $M$ halts. So no
algorithm decides ergodicity for such a family.
\end{enumerate}
\end{corollary}

\begin{proof}
\emph{A map field for $Q$.} Let $Q$ and $\gamma$ be as in (i), and put
\[
Q'(s\mid y)=\frac{Q(s\mid y)-(1-\gamma)\mathbf 1[s=f_M(y)]}{\gamma}.
\]
By hypothesis $Q'(\cdot\mid y)$ is a probability law on $\A_M$ for every $y$.
Enumerate $\A_M=\{s_1,\dots,s_n\}$. For $y\in\A_M^{\mathcal N}$ and $w\in(0,1]$ let
$g(y,w)=s_r$ for the least $r$ with $\sum_{r'\le r}Q'(s_{r'}\mid y)\ge w$. If $W$ is
uniform on $(0,1]$, then $g(y,W)$ has law $Q'(\cdot\mid y)$. Let $\chi_t(v)$ be
Bernoulli with parameter $\gamma$ and $W_t(v)$ uniform on $(0,1]$, all independent.
Put $f_{t,v}=f_M$ if $\chi_t(v)=0$ and $f_{t,v}=g(\cdot,W_t(v))$ if $\chi_t(v)=1$.
This is a local map field, and it is $\gamma$-bounded because
$\Prob(\chi=1\text{ on }S)=\gamma^{|S|}$.

Let $X$ be its run from $x$. Then $X_0,\dots,X_t$ are functions of $x$ and of the
variables of times at most $t$. The variables $(\chi_{t+1}(v),W_{t+1}(v))$, over all
$v$, are independent of these and of each other. So, given $X_0,\dots,X_t$, the
values $X_{t+1}(v)$ are independent, and
\[
\Prob\bigl(X_{t+1}(v)=s\bigm|X_0,\dots,X_t\bigr)
=(1-\gamma)\mathbf 1\bigl[s=f_M((X_t)_v)\bigr]+\gamma\,Q'\bigl(s\mid(X_t)_v\bigr)
=Q\bigl(s\mid(X_t)_v\bigr).
\]
Hence the run from $x$ has law $\mathrm{Law}^Q_x$, and the common-map coupling of the
runs from $x$ and from $\bot$ is a coupling of $\mathrm{Law}^Q_x$ and
$\mathrm{Law}^Q_\bot$.

\emph{Constants.} Recall $\eps_0=(4D_0^2)^{-a}$, with $q\ge3$, $D_0\ge150$ and
$a\ge19$ (\cref{pf:threshold}). From $\gamma\le2\eps_0$ we get
$\beta=2D_0^2\gamma^{1/a}\le2^{1/a}/2\le2^{-18/19}<0.52$. Since $a/q^2>2$ and
$4D_0^2>1$,
\[
\gamma^{1/q^2}\le2^{1/q^2}(4D_0^2)^{-2},\qquad\text{so}\qquad
\lambda=2D_0^2\gamma^{1/q^2}\le\frac{2^{1/9}}{8D_0^2}<10^{-5},
\]
and $2\lambda/(1-\lambda)<0.01$. Also $\gamma\le2\eps_0<10^{-8}$, so
$1-2.02\gamma>0$.

\emph{(i), $M$ halts.} By \cref{thm:width}, or \cref{thm:three} with the constants
of \cref{cor:eight}, $R_M$ is $j$-token refutable with time $\tau=\tau(R_M)$, and
$k\ge\max(2,2j)$. For a window $O$ and a start $x$, \cref{pf:field}(ii) with $y=\bot$
gives
\[
\dtv\bigl(\mathrm{Law}^Q_x(X_t|_O),\mathrm{Law}^Q_\bot(X_t|_O)\bigr)
\le\sum_{v\in O}\Prob\bigl(X_t(v)\ne Y_t(v)\bigr)
\le2|O|\,\gamma^{-\tau/a}\,\frac{\beta^t}{1-\beta},
\]
which tends to $0$. The process forgets, and it is ergodic by
\cref{set:forget-ergodic}.

\emph{(i), $M$ does not halt.} $R_M$ tiles (\cref{thm:width}). By
\cref{pf:field}(i) and the constants, both probabilities are at most
$\gamma(1+2\lambda/(1-\lambda))<1.01\gamma$. As in the proof of \cref{thm:main},
$\Prob^Q_\xi(X_t(v)=\xi(v))>1-1.01\gamma$ and
$\Prob^Q_\bot(X_t(v)=\xi(v))<1.01\gamma$ for every $t$. So the process has memory
with $\alpha=1-2.02\gamma>0$, and \cref{set:memory-nonergodic} shows that it is not
ergodic.

(ii) For every $y$,
\[
\dtv\bigl(Q(\cdot\mid y),p_{M,\eps}(\cdot\mid y)\bigr)
\ge p_{M,\eps}(f_M(y)\mid y)-Q(f_M(y)\mid y)
=1-\eps+\frac{\eps}{|\A_M|}-Q(f_M(y)\mid y).
\]
So $Q(f_M(y)\mid y)\ge1-\eps-d(Q,p_{M,\eps})\ge1-2\eps_0$. Apply (i) with
$\gamma=2\eps_0$.

(iii) By (i), the set of $M$ with $Q_M$ ergodic is the set of halting machines,
which is not decidable.
\end{proof}

\begin{remark}[Scope of \cref{cor:robust}]\label{pf:robust-scope}
\begin{enumerate}[label=\textup{(\alph*)}]
\item \emph{Uniformity needs the growing alphabet.} The ball in (ii) lives in the
space of local kernels on $\A_M$, whose dimension grows with $M$. This is
unavoidable. For one fixed finite alphabet and neighbourhood, the local kernels form
a compact metric space under $d$. Suppose the members of a computable family with
rational entries kept their label on closed balls of one fixed rational radius $r$.
Then members with different labels would lie at distance more than $r$. Take a
finite net of rational kernels with mesh less than $r/2$. Classify each member by a
net point at distance less than $r/2$, which exists and is found by exact rational
computation. Two members with the same net point lie within distance $r$ of each
other, so they share a label, and a finite table decides the labels.
\item \emph{Positive rates.} The ball in (ii) contains kernels with zero entries,
among them the noiseless rule $\delta_{f_M}$. Its positive-rate members, among them
$p_{M,\eps}$, are positive-rate probabilistic cellular automata in the sense of
Theorem~A. As in \cref{pf:scope}, their least transition probability is not bounded
below over the family.
\item \emph{The joint bound is used.} What the proof of (i) uses is a
$\gamma$-bounded common-map representation. For a kernel whose sites update
independently given the past, the construction in the proof provides one. Closeness
of the one-site conditional laws alone is not enough. Consider the transition that,
with probability $\gamma$, sets every site to $\bot$ at once, and otherwise updates
the sites independently by $p_{M,\eps}$. Each one-site conditional law is within
$\gamma$ of $p_{M,\eps}$. Couple the runs by sharing the reset coin and all
replacement variables. Then the runs from any two starts agree after the first
reset, so this process forgets and is ergodic for every $M$. It has no
$\gamma'$-bounded common-map representation for any $\gamma'<1$: from $\xi$, a reset
makes every site of any finite set wrong at once, with probability $\gamma$.
\item \emph{Site-dependent kernels.} The proof of (i) builds the map field site by
site. So (i) holds verbatim when each site $v$ has its own kernel $Q_v$ satisfying
the hypothesis with one $\gamma$.
\item \emph{Correlated faults.} Suppose the layers
$(\chi_t(\cdot),f_{t,\cdot})$, $t=1,2,\dots$, of a local map field are independent
and identically distributed, and each layer is $\gamma$-bounded for sets of sites at
one time, with arbitrary correlations inside a layer. Then:
\begin{itemize}
\item the field is $\gamma$-bounded, by independence of the layers;
\item the runs form a time-homogeneous Markov process in which $X_1|_O$ depends on
$X_0$ only on $O+\mathcal N$ and on randomness independent of $X_0$;
\item so \cref{set:feller,set:forget-ergodic,set:memory-nonergodic} apply verbatim;
\item and with $\gamma\le2\eps_0$, \cref{pf:field} and the constants above give the
dichotomy of (i), with the same bounds, although the process is not a probabilistic
cellular automaton.
\end{itemize}
For a $\gamma$-bounded field that is not time-homogeneous, \cref{pf:field} still
gives the separation of (i) at every time and forgetting under the common-map
coupling, but invariant measures are not defined. Forgetting needs the maps to read
only $\mathcal N$ at the previous time and both runs to use the same map. Memory
needs neither (\cref{pf:field}(i), \cref{mem:laws}).
\end{enumerate}
\end{remark}
